\section{Análisis complementarios}
\label{sec:met_complementarios}

\subsection{Protocolos de inferencia}
\pendiente{Una tabla única indica, para cada resultado reportado, si se usa vista única o TTA, qué
\emph{checkpoint} y qué partición.}

\subsection{Análisis estadístico}
\pendiente{Se indica qué intervalo o prueba corresponde a cada métrica y comparación: \emph{bootstrap}
pareado para diferencias de métricas agregadas, McNemar para aciertos por imagen y Wilcoxon para
\emph{recall} por clase.}

\subsection{Validación externa, perturbaciones, calibración e interpretabilidad}
\pendiente{Además de las métricas sobre el test, se realizan cuatro análisis que responden preguntas
distintas de ``cuánto acierta el modelo'': si el desempeño se sostiene en otra población (evaluación en
HIBA), si se sostiene bajo peores condiciones de captura (perturbaciones de adquisición: ruido, desenfoque,
compresión JPEG y pérdida de resolución), si las probabilidades que entrega el modelo son confiables
(calibración post-hoc ajustada en validación) y si el modelo mira la lesión o el contexto de la imagen
(Grad-CAM). Esta subsección describe el procedimiento de cada uno.}
